\documentclass[reprint,aps,pra,superscriptaddress,nofootinbib,longbibliography]{revtex4-2}

\usepackage{amsmath,amssymb,amsthm}
\usepackage{bm}
\usepackage{mathtools}
\usepackage{booktabs}
\usepackage{microtype}
\usepackage{hyperref}
\hypersetup{hidelinks}

\newtheorem{theorem}{Theorem}
\newtheorem{proposition}[theorem]{Proposition}
\newtheorem{corollary}[theorem]{Corollary}
\newtheorem{definition}[theorem]{Definition}
\newtheorem{lemma}[theorem]{Lemma}
\newtheorem{remark}[theorem]{Remark}

\newcommand{\ket}[1]{|#1\rangle}
\newcommand{\bra}[1]{\langle #1|}
\newcommand{\Tr}{\operatorname{Tr}}
\newcommand{\Repart}{\operatorname{Re}}
\newcommand{\Impart}{\operatorname{Im}}
\newcommand{\diag}{\operatorname{diag}}
\newcommand{\sech}{\operatorname{sech}}
\newcommand{\argmax}{\operatorname*{arg\,max}}
\newcommand{\sig}{\boldsymbol{\sigma}}
\newcommand{\Id}{\mathbb{I}}
\newcommand{\Veff}{V_{\mathrm{eff}}}
\newcommand{\Ftel}{\mathcal F_{\mathrm{tel}}}
\newcommand{\Nneg}{\mathcal N}
\newcommand{\Ddec}{\mathcal D}
\newcommand{\Kbr}{K_{\mathcal E}}

\begin{document}

\title{Local Balancing and Branch Decoherence in Projected GHZ Resources}

\author{Andrei Preece}
\email{andre.preece@web.de}
\affiliation{Independent Researcher, Berlin, Germany}

\author{Boris Batenin}
\email{batenin@web.de}
\affiliation{Independent Researcher, Berlin, Germany}

\date{May 2026}

\begin{abstract}
Projected GHZ resources produce endpoint states whose usefulness is controlled by filtered branch coherence rather than raw multipartite longitudinal parity.  Here we refine this transfer law into an operational error separation.  Any successful completely positive trace-nonincreasing operation on the intermediate subsystem defines a positive branch-transfer kernel
\[
\Kbr=\begin{pmatrix}A&\Gamma\\ \Gamma^*&B\end{pmatrix}\ge0,
\]
where \(A\) and \(B\) are the filtered branch weights and \(\Gamma\) is the filtered branch coherence.  For the two-branch resource \(c\ket{0}^{\otimes N}+s\ket{1}^{\otimes N}\), the endpoint visibility is
\[
\Veff=\frac{2cs|\Gamma|}{c^2A+s^2B}.
\]
Introducing the coherence quality
\[
q=\frac{|\Gamma|}{\sqrt{AB}}
\]
and the branch-balance ratio
\[
r=\frac{c\sqrt A}{s\sqrt B},
\]
we obtain the exact design factorization
\[
\Veff=q\frac{2r}{1+r^2}=q\,\sech(\ln r).
\]
Thus endpoint degradation separates into a coherence-quality factor \(q\) that cannot be increased by branch-preserving local balancing and a branch-imbalance penalty.  We prove that branch-preserving local filters leave \(q\) invariant whenever both filtered branches remain present, can set \(r=1\) whenever both endpoint branches are present, and achieve the maximal branch-preserving filtered concurrence \(C_{\max}=q\) with conditional balancing success probability \(1-|m|\), where \(m=\langle Z\otimes I\rangle\).  We also give a tomography rule that reconstructs \(\Veff\), \(\arg\Gamma\), \(r\), and \(q\) from endpoint observables.  The results provide a diagnostic calculus for separating correctable branch imbalance from branch decoherence that is not corrected by branch-preserving local balancing in projected GHZ and repeater-motivated resources.  All statements are finite-dimensional quantum-mechanical identities.
\end{abstract}

\maketitle

\section{Introduction}

Projected GHZ resources~\cite{GHZ1989,Mermin1990} provide a clean setting in which multipartite resource correlations and endpoint correlations can be different objects.  Entanglement-swapping and repeater-motivated protocols make this distinction operational~\cite{Zukowski1993,Briegel1998,Kimble2008}.  In a two-branch GHZ-type resource, the intermediate operation can preserve, attenuate, or erase coherence between the two global branches.  The endpoint pair inherits the filtered branch coherence, not the raw longitudinal parity coefficient of the original resource.

The preceding transfer law may be stated in prose as follows: the longitudinal multipartite parity coefficient does not determine endpoint entanglement, while the transferred off-diagonal branch coherence \(\Gamma\) does.  This paper asks the next engineering question: once \(\Gamma\) controls endpoint usefulness, which part of the loss is locally correctable and which part is not?

The answer is an exact factorization.  The endpoint visibility can be written as
\begin{equation}
\Veff=q\,\sech(\ln r),
\label{eq:main-factorization-intro}
\end{equation}
where \(q\in[0,1]\) measures normalized branch coherence and \(r>0\) measures branch imbalance.  The two variables play different operational roles.  The factor \(q\) is invariant under branch-preserving local balancing filters.  The factor \(\sech(\ln r)\) is a balance penalty that can be removed probabilistically when both endpoint branches are present.

Thus the central distinction is
\begin{equation}
\boxed{\begin{gathered}
\text{branch imbalance is locally balanceable;}\\
\text{branch decoherence is not.}
\end{gathered}}
\end{equation}
This is not a new dynamical model.  It is a transfer and filtering calculus for a standard class of two-branch endpoint states.  The same calculation gives endpoint concurrence, negativity, teleportation fidelity, and two-terminal CHSH strength in terms of the same single visibility parameter.

\section{Projected GHZ endpoint state}

Consider the GHZ-type resource
\begin{equation}
\ket{G^{(N)}_\eta}
=
c\ket{0}^{\otimes N}+s\ket{1}^{\otimes N},
\qquad
c=\cos\eta,
\quad
s=\sin\eta,
\label{eq:ghz-resource}
\end{equation}
with endpoint qubits \(A\) and \(B\), and middle subsystem
\begin{equation}
M=\{2,\ldots,N-1\}.
\end{equation}
Write
\begin{equation}
\ket{0_M}=\ket{0}^{\otimes(N-2)},
\qquad
\ket{1_M}=\ket{1}^{\otimes(N-2)}.
\end{equation}
Let \(\mathcal E\) be a completely positive trace-nonincreasing map corresponding to one successful intermediate outcome.  Define
\begin{align}
A&=\Tr\mathcal E(\ket{0_M}\bra{0_M}),
&
B&=\Tr\mathcal E(\ket{1_M}\bra{1_M}),\nonumber\\
\Gamma&=\Tr\mathcal E(\ket{0_M}\bra{1_M}),
&
p&=c^2A+s^2B.
\label{eq:ABGammaP}
\end{align}
We assume \(p>0\).  After the successful operation and tracing out the middle output, the normalized endpoint state is
\begin{equation}
\rho_{AB}
=
\frac{1}{c^2A+s^2B}
\left(
\begin{aligned}
&c^2A\ket{00}\bra{00}
+s^2B\ket{11}\bra{11}\\
&+cs\Gamma\ket{00}\bra{11}
+cs\Gamma^*\ket{11}\bra{00}
\end{aligned}
\right).
\label{eq:endpoint-state}
\end{equation}
This is a maximally correlated two-qubit \(X\)-state supported on
\begin{equation}
\mathrm{span}\{\ket{00},\ket{11}\}.
\end{equation}

\section{Branch-transfer kernel}

\begin{definition}[Branch-transfer kernel]
The successful middle operation \(\mathcal E\) defines the branch-transfer kernel
\begin{equation}
\Kbr=
\begin{pmatrix}
A&\Gamma\\
\Gamma^*&B
\end{pmatrix}.
\label{eq:branch-kernel}
\end{equation}
\end{definition}

\begin{lemma}[Kernel positivity]
For every completely positive trace-nonincreasing map \(\mathcal E\),
\begin{equation}
\Kbr\ge0,
\qquad
|\Gamma|^2\le AB.
\label{eq:kernel-positive}
\end{equation}
\end{lemma}

\begin{proof}
Let \(\mathcal E(X)=\sum_\alpha K_\alpha X K_\alpha^\dagger\) be a Kraus representation~\cite{NielsenChuang}.  Define the branch-image vectors in the direct-sum output space by
\begin{equation}
v_0=\bigoplus_\alpha K_\alpha\ket{0_M},
\qquad
v_1=\bigoplus_\alpha K_\alpha\ket{1_M}.
\end{equation}
Then
\begin{equation}
A=\langle v_0|v_0\rangle,
\qquad
B=\langle v_1|v_1\rangle,
\qquad
\Gamma=\langle v_1|v_0\rangle .
\end{equation}
Therefore
\begin{equation}
\widehat K=
\begin{pmatrix}
A&\Gamma^*\\
\Gamma&B
\end{pmatrix}
\end{equation}
is the Gram matrix of \(v_0,v_1\), so \(\widehat K\ge0\).  The displayed branch-transfer kernel \(\Kbr=\widehat K^*\) is its complex conjugate and is also positive semidefinite.  Hence \(\Kbr\ge0\).  Positivity of a \(2\times2\) matrix gives
\begin{equation}
\det\Kbr=AB-|\Gamma|^2\ge0,
\end{equation}
which proves Eq.~\eqref{eq:kernel-positive}.
\end{proof}

The normalized determinant
\begin{equation}
\Ddec=1-\frac{|\Gamma|^2}{AB}
=\frac{\det\Kbr}{AB}
\label{eq:Ddec}
\end{equation}
will be called the branch-decoherence determinant whenever \(AB>0\).  It vanishes for coherent rank-one branch transfer and equals one when \(\Gamma=0\) and \(AB>0\).

\section{Design factorization}

Define the normalized coherence quality and balance ratio by
\begin{equation}
q=\frac{|\Gamma|}{\sqrt{AB}},
\qquad
r=\frac{c\sqrt A}{s\sqrt B},
\label{eq:q-r-def}
\end{equation}
on the open two-branch sector \(csAB>0\).  If one endpoint branch is absent, positivity gives \(\Gamma=0\), the endpoint visibility is \(\Veff=0\), and branch balancing is impossible.  These one-branch cases are treated separately.

\begin{theorem}[Design factorization]
The endpoint phase visibility, concurrence, and transverse tensor radius are
\begin{equation}
\Veff=\frac{2cs|\Gamma|}{c^2A+s^2B}
=q\frac{2r}{1+r^2}
=q\,\sech(\ln r).
\label{eq:design-factorization}
\end{equation}
Consequently,
\begin{equation}
0\le \Veff\le q\le1,
\label{eq:Vleq}
\end{equation}
with \(\Veff=q\) if and only if \(r=1\), equivalently
\begin{equation}
c^2A=s^2B.
\label{eq:balance-condition}
\end{equation}
\end{theorem}

\begin{proof}
Substituting \(|\Gamma|=q\sqrt{AB}\) into \(\Veff\) gives
\begin{equation}
\Veff=
q\frac{2cs\sqrt{AB}}{c^2A+s^2B}.
\end{equation}
Dividing numerator and denominator by \(s^2B\) gives
\begin{equation}
\Veff=q\frac{2r}{1+r^2}.
\end{equation}
With \(x=\ln r\), one has
\begin{equation}
\frac{2r}{1+r^2}
=\frac{2e^x}{1+e^{2x}}
=\frac{1}{\cosh x}
=\sech x.
\end{equation}
The function \(2r/(1+r^2)\) is at most one, with equality exactly at \(r=1\).
\end{proof}

Equivalently, setting \(r=\tan(\theta_{\rm bal}/2)\) gives
\begin{equation}
\Veff=q\sin\theta_{\rm bal}.
\label{eq:theta-bal}
\end{equation}
Thus \(q\) is a radial coherence-quality variable and \(\theta_{\rm bal}\) is a branch-balance angle.

For small imbalance \(x=\ln r\),
\begin{equation}
\Veff=q\left[1-\frac{x^2}{2}+O(x^4)\right].
\label{eq:small-imbalance}
\end{equation}
Branch imbalance is second order near the optimum, whereas coherence loss enters linearly through \(q\).

\section{Local balancing theorem}

Write the endpoint state as
\begin{equation}
\rho_{AB}
=u\ket{00}\bra{00}+v\ket{11}\bra{11}
+w\ket{00}\bra{11}+w^*\ket{11}\bra{00},
\label{eq:uv-state}
\end{equation}
where
\begin{equation}
u=\frac{c^2A}{p},
\qquad
v=\frac{s^2B}{p},
\qquad
w=\frac{cs\Gamma}{p},
\qquad
u+v=1.
\label{eq:uvw}
\end{equation}
Then
\begin{equation}
\Veff=2|w|,
\qquad
q=\frac{|w|}{\sqrt{uv}},
\qquad
r=\sqrt{\frac{u}{v}}.
\label{eq:q-r-uv}
\end{equation}

\begin{theorem}[Branch-preserving local balancing]
Assume \(u,v>0\).  Among successful diagonal branch-preserving local filters acting on one endpoint with \(\alpha\beta>0\), the normalized coherence quality \(q=|w|/\sqrt{uv}\) is invariant, and the post-filter concurrence satisfies
\begin{equation}
C'(\rho_{AB})\le q.
\label{eq:Cprime-bound}
\end{equation}
Equality holds exactly at branch balance.  A one-sided local filter can reach this optimum, producing
\begin{equation}
u'=v'=\frac12,
\qquad
C(\rho'_{AB})=q.
\label{eq:balanced-state-result}
\end{equation}
The maximal conditional success probability of such balancing is
\begin{equation}
P_{\rm bal}=2\min(u,v)=1-|m|,
\qquad
m=\langle Z\otimes I\rangle=u-v.
\label{eq:Pbal}
\end{equation}
\end{theorem}

\begin{proof}
Consider a one-sided diagonal filter
\begin{equation}
F=\alpha\ket{0}\bra{0}+\beta\ket{1}\bra{1},
\qquad
0<\alpha,\beta\le1.
\end{equation}
Applied to endpoint \(A\), it transforms
\begin{equation}
u\mapsto \tilde u=\alpha^2u,
\qquad
v\mapsto \tilde v=\beta^2v,
\qquad
w\mapsto \tilde w=\alpha\beta w.
\end{equation}
After normalization by \(P=\tilde u+\tilde v\), the ratio
\begin{equation}
\frac{|w'|}{\sqrt{u'v'}}
=
\frac{|\tilde w|/P}{\sqrt{(\tilde u/P)(\tilde v/P)}}
=
\frac{|w|}{\sqrt{uv}}
=q
\end{equation}
is invariant.

The concurrence after this successful filter is
\begin{equation}
C'=2|w'|
=
\frac{2\alpha\beta |w|}{\alpha^2u+\beta^2v}
=
q\,
\frac{2\alpha\beta\sqrt{uv}}{\alpha^2u+\beta^2v}.
\end{equation}
Using \(2XY\le X^2+Y^2\) with \(X=\alpha\sqrt u\) and \(Y=\beta\sqrt v\), one obtains
\begin{equation}
C'\le q,
\end{equation}
with equality if and only if
\begin{equation}
\alpha^2u=\beta^2v.
\label{eq:filter-balance-condition}
\end{equation}
To balance the populations, choose \(\alpha^2u=\beta^2v\).  If \(u\ge v\), take \(\beta=1\) and \(\alpha=\sqrt{v/u}\).  If \(v\ge u\), take \(\alpha=1\) and \(\beta=\sqrt{u/v}\).  The conditional balancing success probability is then
\begin{equation}
P_{\rm bal}=2\min(u,v).
\end{equation}
The balanced state has \(|w'|=q/2\), so its concurrence is \(C(\rho'_{AB})=2|w'|=q\).  Since \(m=u-v\) and \(u+v=1\), one has \(2\min(u,v)=1-|m|\).
\end{proof}

\begin{corollary}[Correctable and branch-preserving non-increasable parts]
Branch-preserving local filtering can remove the imbalance penalty \(\sech(\ln r)\), but it cannot increase \(q\).  Thus \(q\) is not increased by branch-preserving local balancing, while \(r\) is the correctable branch-balance factor.
\end{corollary}

The theorem is a Procrustean concentration statement~\cite{Bennett1996} specialized to projected GHZ endpoint states: imbalance can be filtered away probabilistically, but dephased branch coherence cannot be restored by branch-preserving local balancing.

\section{Tomographic separation of errors}

The endpoint tensor in the \(\{x,y,z\}\) basis has transverse entries
\begin{equation}
T^{(2)}_{xx}=2\Repart w,
\qquad
T^{(2)}_{xy}=-2\Impart w,
\qquad
T^{(2)}_{yy}=-2\Repart w.
\end{equation}
Therefore
\begin{equation}
\Veff=
\sqrt{(T^{(2)}_{xx})^2+(T^{(2)}_{xy})^2},
\label{eq:V-tomo}
\end{equation}
and
\begin{equation}
\arg\Gamma=\operatorname{atan2}(-T^{(2)}_{xy},T^{(2)}_{xx}).
\label{eq:phase-tomo}
\end{equation}
The local marginal
\begin{equation}
m=\langle Z\otimes I\rangle=u-v
\label{eq:m-def}
\end{equation}
gives
\begin{equation}
r=\sqrt{\frac{1+m}{1-m}},
\qquad
\frac{2r}{1+r^2}=\sqrt{1-m^2}.
\label{eq:r-from-m}
\end{equation}
Hence, for \(|m|<1\),
\begin{equation}
q=\frac{\Veff}{\sqrt{1-m^2}}.
\label{eq:q-tomo}
\end{equation}

\begin{theorem}[Endpoint tomography of branch quality]
The three endpoint measurements \(T^{(2)}_{xx}\), \(T^{(2)}_{xy}\), and \(m=\langle Z\otimes I\rangle\) determine \(\Veff\) and \(r\) by Eqs.~\eqref{eq:V-tomo}--\eqref{eq:r-from-m}.  On the open two-branch sector \(|m|<1\), they also determine \(q\) by Eq.~\eqref{eq:q-tomo}.  The phase \(\arg\Gamma\) is operationally defined and reconstructed by Eq.~\eqref{eq:phase-tomo} whenever \(\Veff>0\).
\end{theorem}

If the success probability \(p\) and the preparation parameters \(c,s\), with \(cs>0\), are known, the complex branch coherence itself is reconstructed as
\begin{equation}
\Gamma=\frac{p}{2cs}
\left(T^{(2)}_{xx}-iT^{(2)}_{xy}\right).
\label{eq:Gamma-reconstruct}
\end{equation}
Thus endpoint tomography can distinguish two failure modes: \(q<1\) indicates branch decoherence, while \(r\ne1\) indicates branch imbalance.

\section{Endpoint usefulness}

For the endpoint state \eqref{eq:endpoint-state}, the concurrence~\cite{Wootters1998}, negativity, optimal teleportation fidelity, and two-qubit CHSH strength are all functions of \(\Veff\).

\begin{theorem}[One visibility controls endpoint usefulness]
For the projected endpoint state,
\begin{equation}
C(\rho_{AB})=\Veff,
\qquad
\Nneg(\rho_{AB})=\frac{\Veff}{2},
\label{eq:C-N}
\end{equation}
\begin{equation}
\Ftel^{\max}=\frac{2+\Veff}{3},
\label{eq:Ftel}
\end{equation}
and
\begin{equation}
S_{\max}^{(2)}=2\sqrt{1+\Veff^2}.
\label{eq:CHSH}
\end{equation}
\end{theorem}

\begin{proof}
The state is an \(X\)-state with the only off-diagonal coherence \(w\), so its concurrence is \(C=2|w|=\Veff\).  Its partial transpose has eigenvalues \(u,v,|w|,-|w|\), giving \(\Nneg=|w|=\Veff/2\).  The fully entangled fraction is
\begin{equation}
f_{\max}=\max_\phi
\bra{\Phi_\phi}\rho_{AB}\ket{\Phi_\phi}
=\frac12+|w|=\frac{1+\Veff}{2},
\end{equation}
where \(\ket{\Phi_\phi}=(\ket{00}+e^{i\phi}\ket{11})/\sqrt2\).  The standard teleportation fidelity is \(\Ftel^{\max}=(2f_{\max}+1)/3=(2+\Veff)/3\) \cite{Bennett1993,Horodecki1999}.  Finally, the transverse block of the two-qubit correlation tensor is
\begin{equation}
2\begin{pmatrix}
\Repart w&-\Impart w\\
-\Impart w&-\Repart w
\end{pmatrix},
\end{equation}
and hence
\begin{equation}
\left(T^{(2)}\right)^T T^{(2)}
=\diag(\Veff^2,\Veff^2,1).
\end{equation}
The singular values are therefore \(\{1,\Veff,\Veff\}\), and the Horodecki criterion gives Eq.~\eqref{eq:CHSH} \cite{Horodecki1995}.
\end{proof}

Combining this theorem with Eq.~\eqref{eq:design-factorization} gives
\begin{align}
C&=q\sech(\ln r),\nonumber\\
\Ftel^{\max}&=\frac{2+q\sech(\ln r)}{3},\nonumber\\
S_{\max}^{(2)}&=2\sqrt{1+q^2\sech^2(\ln r)}.
\end{align}

\section{Quality--throughput identities}

The pre-balancing endpoint visibility satisfies the exact identity
\begin{equation}
p\Veff=2cs|\Gamma|.
\label{eq:quality-throughput}
\end{equation}
This product is the unnormalized branch coherence transferred by the successful middle operation.

After optimal branch-preserving local balancing, the additional conditional balancing success probability is \(P_{\rm bal}=1-|m|\), and the final concurrence is \(q\).  The total success probability of successful middle operation followed by balancing is
\begin{equation}
p_{\rm tot}=p(1-|m|)=2\min(c^2A,s^2B),
\label{eq:ptot}
\end{equation}
and its concurrence yield is
\begin{equation}
p_{\rm tot}q
=2\min(c^2A,s^2B)\frac{|\Gamma|}{\sqrt{AB}}.
\label{eq:balanced-yield}
\end{equation}
Thus a protocol can choose between immediate endpoint use, quantified by \(p\Veff\), and a balanced endpoint resource, quantified by \(p_{\rm tot}q\).  These are different operational objectives and need not be optimized by the same filter.

\section{Examples and diagnostics}

\subsection{Balanced coherent projection}

If \(A=B\), \(|\Gamma|=\sqrt{AB}\), and \(c=s=1/\sqrt2\), then
\begin{equation}
q=1,
\qquad
r=1,
\qquad
\Veff=1.
\end{equation}
The endpoint state is Bell-equivalent and reaches
\begin{equation}
C=1,
\qquad
\Ftel^{\max}=1,
\qquad
S_{\max}^{(2)}=2\sqrt2.
\end{equation}

\subsection{Pure imbalance without branch decoherence}

If \(q=1\) but \(r\ne1\), then the endpoint state is a partially entangled pure state with
\begin{equation}
\Veff=\sech(\ln r)<1.
\end{equation}
Local balancing succeeds with conditional probability \(1-|m|\) and produces a Bell state.  The loss is therefore not irreversible decoherence; it is branch imbalance.

\subsection{Pure branch decoherence}

If \(r=1\) but \(q<1\), then
\begin{equation}
\Veff=q.
\end{equation}
No branch-preserving local balancing filter can increase \(q\).  The loss is not corrected within this filtering class.

\subsection{Complete coherence loss}

If \(\Gamma=0\), then
\begin{equation}
q=0,
\qquad
\Veff=0,
\qquad
C=0,
\qquad
S_{\max}^{(2)}=2.
\end{equation}
Nevertheless, before endpoint bit-flip feedforward,
\begin{equation}
T^{(2)}_{zz}=1.
\end{equation}
This shows again that longitudinal endpoint saturation is not an entanglement witness for this class.  If a known endpoint bit-flip feedforward is applied, the endpoint support is mapped from \(\mathrm{span}\{\ket{00},\ket{11}\}\) to \(\mathrm{span}\{\ket{01},\ket{10}\}\), and the endpoint longitudinal entry changes from \(T^{(2)}_{zz}=1\) to \(T^{(2)}_{zz}=-1\).  This feedforward sign convention still does not identify the endpoint longitudinal entry with the resource coefficient \(d_N\).

\section{Scope and limitations}

The results apply to endpoint states generated from two-branch GHZ-type resources by successful completely positive trace-nonincreasing operations on the intermediate subsystem.  They also apply to any two-qubit endpoint state of the maximally correlated form \eqref{eq:uv-state}.

The local balancing theorem is stated for branch-preserving diagonal local filters.  It does not claim that arbitrary multipartite resources, W states, or general mixed states admit a single branch-transfer kernel of the form \eqref{eq:branch-kernel}.  The theory separates mechanisms inside the two-branch sector; it is not a universal theory of all multipartite entanglement.

No hidden variables, modified Born rule, new Bell inequality, or new dynamical model are introduced.  The contribution is an operational decomposition of endpoint usefulness into a correctable branch-balance factor and a branch-decoherence invariant.

\section{Conclusion}

Projected GHZ resources admit an exact operational separation of endpoint degradation.  The branch-transfer kernel
\begin{equation}
\Kbr=
\begin{pmatrix}
A&\Gamma\\
\Gamma^*&B
\end{pmatrix}
\end{equation}
contains two independent ingredients: the normalized branch coherence
\begin{equation}
q=\frac{|\Gamma|}{\sqrt{AB}}
\end{equation}
and the branch-balance ratio
\begin{equation}
r=\frac{c\sqrt A}{s\sqrt B}.
\end{equation}
The endpoint visibility factorizes as
\begin{equation}
\Veff=q\frac{2r}{1+r^2}=q\sech(\ln r).
\end{equation}
Branch-preserving local filters can set \(r=1\) and reach \(C=q\) with conditional balancing success probability \(1-|m|\), but they cannot increase \(q\).  Thus branch imbalance is locally balanceable, while branch decoherence is not corrected within the class of branch-preserving local balancing operations.

Endpoint tomography separates the two mechanisms directly:
\begin{align}
\Veff&=\sqrt{(T^{(2)}_{xx})^2+(T^{(2)}_{xy})^2},\nonumber\\
r&=\sqrt{\frac{1+m}{1-m}},\nonumber\\
q&=\frac{\Veff}{\sqrt{1-m^2}}.
\end{align}
This gives a practical diagnostic rule for projected GHZ and repeater-motivated resources: if \(r\ne1\), rebalance the branches; if \(q<1\), the loss is branch decoherence.  The same visibility determines concurrence, negativity, teleportation fidelity, and CHSH strength.

\begin{acknowledgments}
No external funding was received.
\end{acknowledgments}

\section*{Data Availability}
All calculations are analytic.  A minimal numerical verification script is included with the submission files.

\section*{Author Contributions}
Both authors contributed to the conceptual formulation, analytic derivations, and manuscript preparation.

\begin{thebibliography}{99}

\bibitem{Briegel1998}
H.-J. Briegel, W. D\"ur, J. I. Cirac, and P. Zoller,
Quantum repeaters: The role of imperfect local operations in quantum communication,
Phys. Rev. Lett. \textbf{81}, 5932 (1998).

\bibitem{Kimble2008}
H. J. Kimble,
The quantum internet,
Nature \textbf{453}, 1023 (2008).

\bibitem{GHZ1989}
D. M. Greenberger, M. A. Horne, and A. Zeilinger,
Going beyond Bell's theorem,
in \textit{Bell's Theorem, Quantum Theory and Conceptions of the Universe},
edited by M. Kafatos (Kluwer, Dordrecht, 1989), pp. 69--72.

\bibitem{Mermin1990}
N. D. Mermin,
Extreme quantum entanglement in a superposition of macroscopically distinct states,
Phys. Rev. Lett. \textbf{65}, 1838 (1990).

\bibitem{Zukowski1993}
M. \.{Z}ukowski, A. Zeilinger, M. A. Horne, and A. K. Ekert,
``Event-ready-detectors'' Bell experiment via entanglement swapping,
Phys. Rev. Lett. \textbf{71}, 4287 (1993).

\bibitem{Bennett1993}
C. H. Bennett, G. Brassard, C. Cr\'epeau, R. Jozsa, A. Peres, and W. K. Wootters,
Teleporting an unknown quantum state via dual classical and Einstein-Podolsky-Rosen channels,
Phys. Rev. Lett. \textbf{70}, 1895 (1993).

\bibitem{Bennett1996}
C. H. Bennett, H. J. Bernstein, S. Popescu, and B. Schumacher,
Concentrating partial entanglement by local operations,
Phys. Rev. A \textbf{53}, 2046 (1996).

\bibitem{Horodecki1995}
R. Horodecki, P. Horodecki, and M. Horodecki,
Violating Bell inequality by mixed spin-\(\frac12\) states: necessary and sufficient condition,
Phys. Lett. A \textbf{200}, 340 (1995).

\bibitem{Horodecki1999}
M. Horodecki, P. Horodecki, and R. Horodecki,
General teleportation channel, singlet fraction, and quasidistillation,
Phys. Rev. A \textbf{60}, 1888 (1999).

\bibitem{Wootters1998}
W. K. Wootters,
Entanglement of formation of an arbitrary state of two qubits,
Phys. Rev. Lett. \textbf{80}, 2245 (1998).

\bibitem{NielsenChuang}
M. A. Nielsen and I. L. Chuang,
\textit{Quantum Computation and Quantum Information}
(Cambridge University Press, Cambridge, 2000).

\end{thebibliography}

\end{document}
